\begin{proof}[Proof of \cref{obj:lem:predecessor-score-homeomorphism}]
Write \(B=B_i^\tau\), \(\mathcal Z=\mathcal Z_i^\tau\), and \(E=E_i^\tau\). For \(\chi\in\mathcal Z\), set \(v_\chi=E(\chi)\), so
\[
(v_\chi)_a
=
\begin{cases}
\chi_b, & a=\pi(b)\text{ for some }b\in B,\\
0, & a\notin \pi(B).
\end{cases}
\]
This point lies in \([0,1]^n\), since each displayed predecessor coordinate lies in \([0,1]\) and every filled coordinate is \(0\).

\begin{enumerate}
\item The map \(\Lambda_i^\tau\) is continuous. Indeed, \(E:\mathcal Z\to[0,1]^n\) is continuous coordinatewise: each coordinate is either one of the product-coordinate projections \(\chi\mapsto\chi_b\) or the constant \(0\). For each \(b\in B\), positivity gives
\[
p_{\pi(b)}\!\left((v_\chi)_{\pi(b)}\mid (v_\chi)_{\operatorname{pa}_G(\pi(b))}\right)>0,
\qquad
q_{\pi(b)}\!\left((v_\chi)_{\pi(b)}\right)>0,
\]
and the \(C^3\) regularity in the hypotheses gives continuity of both factors on the relevant closed cubes. Hence the quotient is a continuous positive function of \(\chi\), and applying \(\log\) preserves continuity. Since this holds for every coordinate \(b\in B\), the product map \(\Lambda_i^\tau:\mathcal Z\to\mathbb R^B\) is continuous. % lean: continuous_predecessorScoreMap

\item The map \(\Lambda_i^\tau\) is injective. Let \(\chi,\chi'\in\mathcal Z\) satisfy
\[
\Lambda_i^\tau(\chi)=\Lambda_i^\tau(\chi'),
\]
and put \(v=v_\chi\), \(w=v_{\chi'}\). With \(L_b=\log R_b\), the latent ratio identity in the perfect-intervention hypotheses gives, for every \(b\in B\),
\[
L_b(f(v))
=
\log
\frac{q_{\pi(b)}(v_{\pi(b)})}
{p_{\pi(b)}(v_{\pi(b)}\mid v_{\operatorname{pa}_G(\pi(b))})},
\qquad
L_b(f(w))
=
\log
\frac{q_{\pi(b)}(w_{\pi(b)})}
{p_{\pi(b)}(w_{\pi(b)}\mid w_{\operatorname{pa}_G(\pi(b))})}.
\]
Thus equality of the score vectors is exactly equality of the predecessor log-ratio projections:
\[
\bigl(L_b(f(v))\bigr)_{b\in B}
=
\bigl(L_b(f(w))\bigr)_{b\in B}.
\]

We prove by strong induction on \(m\in\mathbb N\) that
\[
\forall b\in B,\quad \tau(b)=m\Longrightarrow v_{\pi(b)}=w_{\pi(b)}.
\]
Fix \(b\in B\) with \(\tau(b)=m\), and assume the claim for all smaller order values. If \(a\in\operatorname{pa}_G(\pi(b))\), let \(c=\pi^{-1}(a)\). Then \(\pi(c)\to\pi(b)\) is an edge of \(G\), so the transported latent ordering gives
\[
\tau(c)<\tau(b)=m.
\]
Since \(b\in B\), also \(\tau(b)<\tau(i)\), hence \(\tau(c)<\tau(i)\) and \(c\in B\). The induction hypothesis gives
\[
v_a=v_{\pi(c)}=w_{\pi(c)}=w_a.
\]
Therefore
\[
v_a=w_a
\qquad
\text{for all }a\in\operatorname{pa}_G(\pi(b)).
\]

For \(\zeta\in[0,1]\), let \(v^{[b:\zeta]}\in[0,1]^n\) be the point obtained from \(v\) by replacing the coordinate \(\pi(b)\) by \(\zeta\):
\[
\bigl(v^{[b:\zeta]}\bigr)_a
=
\begin{cases}
\zeta, & a=\pi(b),\\
v_a, & a\ne \pi(b).
\end{cases}
\]
Define, on \([0,1]\), the one-coordinate score
\[
\psi_b(\zeta)
=
\log
\frac{q_{\pi(b)}(\zeta)}
{p_{\pi(b)}\!\left(\zeta\mid v_{\operatorname{pa}_G(\pi(b))}\right)}.
\]
Strict positivity and \(C^3\) regularity make \(\psi_b\) continuous on \([0,1]\) and differentiable on \((0,1)\). For \(0<\zeta<1\), the point \(v^{[b:\zeta]}\) lies in the latent cube, so the fixed own-derivative sign hypothesis gives
\[
0
<
s_{\pi(b)}\,
\frac{d}{d\zeta}
\log
\frac{q_{\pi(b)}(\zeta)}
{p_{\pi(b)}\!\left(\zeta\mid v_{\operatorname{pa}_G(\pi(b))}\right)}
=
s_{\pi(b)}\,\psi_b'(\zeta).
\]
Since \(s_{\pi(b)}\in\{-1,1\}\), the mean-value theorem yields, for any \(0\le \zeta_1<\zeta_2\le1\),
\[
\psi_b(\zeta_2)-\psi_b(\zeta_1)
=
\psi_b'(\zeta_\circ)(\zeta_2-\zeta_1)
\]
for some \(\zeta_\circ\in(\zeta_1,\zeta_2)\), with positive right-hand side when \(s_{\pi(b)}=1\) and negative right-hand side when \(s_{\pi(b)}=-1\). Thus \(\psi_b\) is strictly monotone on \([0,1]\).

Moreover,
\[
\psi_b(v_{\pi(b)})
=
\log
\frac{q_{\pi(b)}(v_{\pi(b)})}
{p_{\pi(b)}(v_{\pi(b)}\mid v_{\operatorname{pa}_G(\pi(b))})},
\]
and parent locality, together with \(v_a=w_a\) for all \(a\in\operatorname{pa}_G(\pi(b))\), gives
\[
p_{\pi(b)}(w_{\pi(b)}\mid v_{\operatorname{pa}_G(\pi(b))})
=
p_{\pi(b)}(w_{\pi(b)}\mid w_{\operatorname{pa}_G(\pi(b))}),
\]
so
\[
\psi_b(w_{\pi(b)})
=
\log
\frac{q_{\pi(b)}(w_{\pi(b)})}
{p_{\pi(b)}(w_{\pi(b)}\mid w_{\operatorname{pa}_G(\pi(b))})}.
\]
The equality \(L_b(f(v))=L_b(f(w))\) therefore gives
\[
\psi_b(v_{\pi(b)})=\psi_b(w_{\pi(b)}).
\]
By strict monotonicity of \(\psi_b\),
\[
v_{\pi(b)}=w_{\pi(b)}.
\]
The induction is complete. Finally, for each \(b\in B\),
\[
\chi_b=v_{\pi(b)}=w_{\pi(b)}=\chi'_b,
\]
so \(\chi=\chi'\) in the product cube. Thus \(\Lambda_i^\tau\) is injective. % lean: predecessorScoreMap_injective

\item Since \(B\) is finite, \(\mathcal Z=\prod_{b\in B}[0,1]\) is compact, and \(\mathbb R^B\) is Hausdorff. Every closed subset of \(\mathcal Z\) is compact, so its image under the continuous map \(\Lambda_i^\tau\) is compact and hence closed in \(\mathbb R^B\). Together with injectivity, this makes \(\Lambda_i^\tau:\mathcal Z\to\mathbb R^B\) a topological embedding with closed image. Therefore the induced map
\[
\Lambda_i^\tau:\mathcal Z\longrightarrow \Lambda_i^\tau(\mathcal Z)=\mathcal R_i^\tau
\]
is a homeomorphism onto its range. Consequently \((\Lambda_i^\tau)^{-1}:\mathcal R_i^\tau\to\mathcal Z_i^\tau\) is continuous. The same argument covers the case \(B=\varnothing\), where both products are singleton spaces. % lean: predecessorScoreHomeomorph
\end{enumerate}
\end{proof}
